\heading{实验物理中的统计方法-矩方法}





\begin{question}
考虑服从高斯分布的随机变量 \(x\)，均值 \(\mu\) 和方差 \(\sigma^2\)
未知，并假设样本为 \(x_1,\ldots,x_n\)。

\begin{enumerate}
\item
  利用矩方法构造 \(\mu\) 和 \(\sigma^2\) 的估计量。利用函数
  \(a_1=x\)，\(a_2=x^2\)，使得期待值 \(E[a_i(x)]\) 对应于 \(x\)
  的一阶和二阶代数矩。
\item
  计算 (a) 中得到的估计量 \(\hat\mu\) 与 \(\hat\sigma^2\)
  的期待值。这两个估计量是否是无偏的？
\end{enumerate}
\begin{solution}
设样本 \(x_1,\dots,x_n\) 来自高斯分布
\(N(\mu,\sigma^2)\)，但这里我们不用最大似然，而用矩方法。

\begin{enumerate}
\item 构造 \(\mu\) 与 \(\sigma^2\) 的矩估计

取

\[
a_1(x)=x,
\qquad
a_2(x)=x^2.
\]

理论矩为

\[
E[x]=\mu,
\qquad E[x^2]=\mu^2+\sigma^2.
\]

样本矩为

\[
\overline x=\frac1n\sum_{i=1}^n x_i,
\qquad
\overline{x^2}=\frac1n\sum_{i=1}^n x_i^2.
\]

令样本矩等于理论矩：

\[
\hat\mu=\overline x,
\qquad
\hat\sigma^2=\overline{x^2}-\overline x^{\,2}.
\]

也就是

\[
\hat\sigma^2=\frac1n\sum_{i=1}^n (x_i-\overline x)^2.
\]

这正好与高斯模型下的方差 MLE 一致。

\item 期望值：是否无偏

显然

\[
E[\hat\mu]=E[\overline x]=\mu,
\]

因此 \(\hat\mu\) 无偏。

对方差估计量，利用第五章结论

\[
s^2=\frac1{n-1}\sum_{i=1}^n (x_i-\overline x)^2
\]

是 \(\sigma^2\) 的无偏估计，而

\[
\hat\sigma^2=\frac{n-1}{n}s^2,
\]

所以

\[
E[\hat\sigma^2]=\frac{n-1}{n}\sigma^2.
\]

因此矩方法得到的 \(\hat\sigma^2\) 不是无偏估计，它向下偏。
\end{enumerate}
\end{solution}
\end{question}



\begin{question}
考虑粒子反应中产生的 \(\rho^0\) 介子衰变为两个带电 \(\pi\)
介子（\(\pi^+\pi^-\)）。衰变角定义为，在 \(\pi^+\pi^-\) 质心系中
\(\pi^+\) 运动方向与 \(\rho\) 的原初方向的夹角，见图 8.1。由于
\(\rho^0\) 的自旋为 1，\(\pi\) 介子的自旋为 0，可以证明 \(\cos\theta\)
的分布具有如下形式
\pyfig[0.62\linewidth]{figures/fig_08_01_rho_decay_schematic.png}{图 8.1：\(\rho^0\to\pi^+\pi^-\) 衰变角定义示意图。}


\[
f(\cos\theta;\eta)=\frac12(1-\eta)+\frac32\eta\cos^2\theta,
\]

其中自旋排列参数 \(\eta\) 的取值范围为 \(-\frac12\le\eta\le1\)。

\begin{enumerate}
\item
  假设某反应产生的 \(\rho^0\) 中，测量到 \(n\) 个 \(\cos\theta\)
  值。利用矩方法，取 \(a=x^2\)，构造自旋排列参数的估计量
  \(\hat\eta\)。为什么不能用 \(a=x\) 构造估计量？
\item
  计算 \(\hat\eta\) 的期待值和方差。
\end{enumerate}
\begin{solution}
定义 \(x=\cos\theta\)，分布为

\[
f(x;\eta)=\frac12(1-\eta)+\frac32\eta x^2,
\qquad -1\le x\le 1,
\qquad -\frac12\le\eta\le 1.
\]

\begin{enumerate}
\item 用矩方法构造 \(\eta\) 的估计量；为何不能取 \(a=x\)

先计算一阶矩：

\[
E[x]=\int_{-1}^1 x f(x;\eta)\,dx = 0.
\]

这与 \(\eta\) 无关，所以用 \(a=x\) 不能区分不同的
\(\eta\)，自然无法构造估计量。

于是取二阶矩 \(a=x^2\)。先算

\[
E[x^2]=\int_{-1}^1 x^2 f(x;\eta)\,dx
=\frac{1-\eta}{2}\int_{-1}^1 x^2dx + \frac{3\eta}{2}\int_{-1}^1 x^4dx.
\]

利用

\[
\int_{-1}^1 x^2dx=\frac23,
\qquad
\int_{-1}^1 x^4dx=\frac25,
\]

得到

\[
E[x^2]=\frac{1-\eta}{3}+\frac{3\eta}{5}=\frac{5+4\eta}{15}.
\]

令样本二阶矩

\[
\overline{x^2}=\frac1n\sum_{i=1}^n x_i^2
\]

等于理论矩，解出

\[
\hat\eta = \frac{15\overline{x^2}-5}{4}.
\]

\item 计算 \(E[\hat\eta]\) 与 \(V[\hat\eta]\)

由构造直接知道

\[
E[\hat\eta]=\eta,
\]

因此它是无偏估计量。

接着要算方差。先求四阶矩：

\[
E[x^4]=\frac{1-\eta}{2}\int_{-1}^1 x^4dx + \frac{3\eta}{2}\int_{-1}^1 x^6dx.
\]

再用

\[
\int_{-1}^1 x^6dx=\frac27,
\]

得到

\[
E[x^4]=\frac{1-\eta}{5}+\frac{3\eta}{7}=\frac{7+8\eta}{35}.
\]

于是

\[
V(\overline{x^2})=\frac1n\left(E[x^4]-E[x^2]^2\right).
\]

而

\[
\hat\eta = \frac{15}{4}\overline{x^2}-\frac54,
\]

故

\[
V[\hat\eta]=\left(\frac{15}{4}\right)^2 V(\overline{x^2}).
\]
\pyfig[0.72\linewidth]{figures/fig_08_02_eta_density.png}{不同 \(\eta\) 下 \(x=\cos\theta\) 的理论密度；对称性说明一阶矩恒为 0。}
代入并化简后得到

\[
V[\hat\eta]=\frac{-28\eta^2+20\eta+35}{28n}.
\]
\end{enumerate}
\end{solution}
\end{question}

